% sections/introduction.tex
\section{Purpose and Scope}

This protocol defines the standard colony-PCR workflow used by the Molecular Cloning Core to
verify insert presence in \textit{Escherichia coli} DH5$\alpha$ colonies following transformation
with the pNX-303 expression construct. The assay distinguishes recombinant colonies carrying the
612\,bp insert-junction fragment from empty-vector religants, which instead yield a 348\,bp
product from the unoccupied multiple cloning site. It applies to all transformation screens
processed under Nexa Biosciences cloning workflows NX-CL-01 through NX-CL-07 and is executed by
trained laboratory staff following the equipment qualification described in NX-EQ-014.

This document serves a dual purpose: it is both the controlled procedure (Sections~2--5) and the
executed record for Run~TC-2026-0729-C (Sections~5--7), completed contemporaneously by the
operator and countersigned by the reviewer in Section~7.

\subsection{Principle}

The polymerase chain reaction amplifies a defined DNA segment through repeated cycles of thermal
denaturation, primer annealing, and polymerase-mediated extension, first described by
\citet{mullis1987specific} and established as a routine molecular biology tool following
\citet{saiki1988primer}. Two oligonucleotide primers (NX-303-F and NX-303-R) flank the
insert--vector junction; their relative positions produce a diagnostic size shift between
recombinant and non-recombinant templates, allowing colony genotype to be read directly from an
agarose gel without plasmid extraction. Reaction assembly and thermocycling parameters below
follow the general practice outlined in \citet{green2019pcr} and \citet{sambrook2001molecular},
adapted to the Nimbus GT-96 thermocycler and the Vertex Taq 2X Master Mix used in this facility.

\begin{notebox}
Colony PCR uses whole cells (or a boiled lysate) directly as template. Cell debris and residual
LB medium are tolerated by the Vertex Taq 2X Master Mix at the input volumes specified in
Section~3, but template volume must not exceed 15\% of the final reaction volume or amplification
efficiency degrades measurably.
\end{notebox}
